\documentclass{beamer}
\usepackage[utf8]{inputenc}

\usetheme{Madrid}
\usecolortheme{default}
\usepackage{amsmath,amssymb,amsfonts,amsthm}
\usepackage{txfonts}
\usepackage{tkz-euclide}
\usepackage{listings}
\usepackage{adjustbox}
\usepackage{array}
\usepackage{tabularx}
\usepackage{gvv}
\usepackage{lmodern}
\usepackage{circuitikz}
\usepackage{tikz}
\usepackage{graphicx}

\setbeamertemplate{page number in head/foot}[totalframenumber]

\usepackage{tcolorbox}
\tcbuselibrary{minted,breakable,xparse,skins}



\definecolor{bg}{gray}{0.95}
\DeclareTCBListing{mintedbox}{O{}m!O{}}{%
	breakable=true,
	listing engine=minted,
	listing only,
	minted language=#2,
	minted style=default,
	minted options={%
		linenos,
		gobble=0,
		breaklines=true,
		breakafter=,,
		fontsize=\small,
		numbersep=8pt,
		#1},
	boxsep=0pt,
	left skip=0pt,
	right skip=0pt,
	left=25pt,
	right=0pt,
	top=3pt,
	bottom=3pt,
	arc=5pt,
	leftrule=0pt,
	rightrule=0pt,
	bottomrule=2pt,
	toprule=2pt,
	colback=bg,
	colframe=orange!70,
	enhanced,
	overlay={%
		\begin{tcbclipinterior}
			\fill[orange!20!white] (frame.south west) rectangle ([xshift=20pt]frame.north west);
	\end{tcbclipinterior}},
	#3,
}
\lstset{
	language=C,
	basicstyle=\ttfamily\small,
	keywordstyle=\color{blue},
	stringstyle=\color{orange},
	commentstyle=\color{green!60!black},
	numbers=left,
	numberstyle=\tiny\color{gray},
	breaklines=true,
	showstringspaces=false,
}
%------------------------------------------------------------
%This block of code defines the information to appear in the
%Title page
\title %optional
{12.476}
%\subtitle{A short story}

\author % (optional)
{RAVULA SHASHANK REDDY - EE25BTECH11047}

 \begin{document}
	
	
	\frame{\titlepage}
	\begin{frame}{Question}
    If
\begin{align*}
\vec{A} \vec{B} = \vec{I},
\end{align*}

then:

\begin{enumerate}[label=\alph*)]
    \item $\vec{B} = \vec{A}^{\top}$
    \item $\vec{A} = \vec{B}^{\top}$
    \item $\vec{B} = \vec{A}^{-1}$
    \item $\vec{B} = \vec{A}$
\end{enumerate}

\end{frame}
\begin{frame}{Solution}
    Given:
\begin{align}
 \vec{A} \vec{B} = \vec{I} 
 \end{align}
 
Multiply both sides by  $\vec{A}^{-1}$: 
\begin{align}
&\vec{A}^{-1} (\vec{A} \vec{B}) = \vec{A}^{-1} \vec{I} \\
\Rightarrow \quad & (\vec{A}^{-1} \vec{A}) \vec{B} = \vec{A}^{-1} \\
\Rightarrow \quad & \vec{I} \, \vec{B} = \vec{A}^{-1} \\
\Rightarrow \quad & \vec{B} = \vec{A}^{-1}
\end{align}
\begin{align}
 \boxed{\vec{B} = \vec{A}^{-1}}.
\end{align}
\end{frame}
\begin{frame}[fragile]
\frametitle{C Code}
\begin{lstlisting}
    #include <stdio.h>

#define N 3  // size of the square matrices

// Function to multiply two matrices
void multiply(int A[N][N], int B[N][N], int C[N][N]) {
    for(int i=0; i<N; i++) {
        for(int j=0; j<N; j++) {
            C[i][j] = 0;
            for(int k=0; k<N; k++)
                C[i][j] += A[i][k] * B[k][j];
        }
    }
}

// Function to print a matrix
void printMatrix(int M[N][N]) {
    for(int i=0; i<N; i++) {
    \end{lstlisting}
    \end{frame}
\begin{frame}[fragile]
\frametitle{C Code}
\begin{lstlisting}
        for(int j=0; j<N; j++)
            printf("%d ", M[i][j]);
        printf("\n");
    }
}

int main() {
    // Example matrices
    int A[N][N] = {{1,0,0},{0,1,0},{0,0,1}};
    int B[N][N] = {{1,0,0},{0,1,0},{0,0,1}}; // B = inverse of A (identity in this case)
    int C[N][N];

    multiply(A, B, C);

    printf("Matrix AB:\n");
    printMatrix(C);
    \end{lstlisting}
    \end{frame}
\begin{frame}[fragile]
\frametitle{C Code}
\begin{lstlisting}
    // Check if C is identity
    int isIdentity = 1;
    for(int i=0;i<N;i++)
        for(int j=0;j<N;j++)
            if((i==j && C[i][j]!=1) || (i!=j && C[i][j]!=0))
                isIdentity = 0;

    if(isIdentity)
        printf("Hence, B = A^-1\n");
    else
        printf("B is NOT the inverse of A\n");

    return 0;
}
\end{lstlisting}
\end{frame}
\begin{frame}[fragile]
\frametitle{Python Direct}
\begin{lstlisting}
    import numpy as np

# Example 3x3 matrix A
A = np.array([[2, 1, 1],
              [1, 2, 1],
              [1, 1, 2]])

# Suppose B satisfies AB = I
# Compute B = inverse of A
B = np.linalg.inv(A)

# Verify AB = I
AB = np.dot(A, B)
\end{lstlisting}
\end{frame}
\begin{frame}[fragile]
\frametitle{Python Direct}
\begin{lstlisting}
print("Matrix A:\n", A)
print("\nMatrix B (A^-1):\n", B)
print("\nAB:\n", AB)

# Check if AB is approximately identity
if np.allclose(AB, np.eye(3)):
    print("\nHence, B = A^-1")
else:
    print("\nB is NOT the inverse of A")
\end{lstlisting}
\end{frame}
\begin{frame}[fragile]
\frametitle{Python Shared}
\begin{lstlisting}
    import ctypes
import numpy as np

N = 3

# Load the shared library
lib = ctypes.CDLL('./libmatrix.so')

# Define argument types for verify_inverse
lib.verify_inverse.argtypes = [ctypes.POINTER(ctypes.c_double) for _ in range(2)]
lib.verify_inverse.restype = ctypes.c_int

# Example matrix A
A = np.array([[2.0,1.0,1.0],
              [1.0,2.0,1.0],
              [1.0,1.0,2.0]])
\end{lstlisting}
\end{frame}
\begin{frame}[fragile]
\frametitle{Python Shared}
\begin{lstlisting}
# Compute B = inverse of A using numpy
B = np.linalg.inv(A)

# Convert numpy arrays to ctypes
A_ctypes = A.ctypes.data_as(ctypes.POINTER(ctypes.c_double))
B_ctypes = B.ctypes.data_as(ctypes.POINTER(ctypes.c_double))

# Call C function
result = lib.verify_inverse(A_ctypes, B_ctypes)

if result:
    print("AB = I, hence B = A^-1 (Verified in C via ctypes)")
else:
    print("B is NOT the inverse of A")
\end{lstlisting}
\end{frame}
\end{document}
